\section*{Chapitre \ref{ch:b1:geological-maps} --- Les cartes géologiques}

\begin{solution}{exo:b1:geological-maps:1}
\omterm{def:b1:geological-maps:attitude}{Direction}~: \omterm{def:b1:spherical-astronomy:horizontal}{azimut} des horizontales du plan~; \omterm{def:b1:geological-maps:attitude}{pendage}~: angle entre le plan et
l'horizontale, mesuré perpendiculairement à la \omterm{def:b1:geological-maps:attitude}{direction}~; \omterm{def:b1:geological-maps:attitude}{sens du pendage}~:
\omterm{def:b1:spherical-astronomy:horizontal}{azimut} de la ligne de plus grande pente, orientée vers le bas. Pour $135/40$, la
\omterm{def:b1:geological-maps:attitude}{direction} par la règle de la main droite est $135 - 90 = \ang{045}$. Symbole~: un
trait de \omterm{def:b1:geological-maps:attitude}{direction} orienté N\,45\,E, avec un tiret sur son côté sud-est (vers
\ang{135}) accompagné de 40.
\end{solution}

\begin{solution}{exo:b1:geological-maps:2}
$z = 500 - \tan\ang{30}\,(x\sin\ang{90} + y\cos\ang{90}) = 500 - 0.577\,x$.
En $(1000, 0)$~: $500 - 577 = \qty{-77}{m}$. En $(0, 1000)$~: \qty{500}{m}, le
plan étant horizontal le long de sa \omterm{def:b1:geological-maps:attitude}{direction} nord--sud.
\end{solution}

\begin{solution}{exo:b1:geological-maps:3}
Une trace qui suit une courbe de niveau est celle d'un contact horizontal, à
\qty{350}{m}. Une trace qui traverse en ligne droite collines et vallées est
celle d'un contact vertical.
\end{solution}

\begin{solution}{exo:b1:geological-maps:4}
Le long de la rivière, le cœur de l'\omterm{def:b1:geological-maps:fold}{anticlinal} montre le schiste et le granite
(sous le grès)~; le cœur du \omterm{def:b1:geological-maps:fold}{synclinal} est occupé par la craie. Les unités les
plus anciennes sont les plus basses de la série~: là où les couches se bombent
vers le haut et où l'érosion les entaille le plus profondément, les plus
anciennes apparaissent~; là où elles s'affaissent, les plus jeunes sont
préservées au milieu.
\end{solution}

\begin{solution}{exo:b1:geological-maps:5}
Les deux premiers points sont à la même altitude~: la droite $y = 0$ est une
ligne de \omterm{def:b1:geological-maps:attitude}{direction}, donc la \omterm{def:b1:geological-maps:attitude}{direction} est est--ouest. Le contact est \qty{100}{m}
plus bas \qty{500}{m} plus au nord~: il plonge vers le nord,
$\tan\delta = 100/500 = 0.2$, $\delta \approx \ang{11.3}$~; \omterm{def:b1:geological-maps:attitude}{sens du pendage}
\ang{000}, \omterm{def:b1:geological-maps:attitude}{direction} \ang{270} par la règle de la main droite. Avec le
théorème~: $\overrightarrow{P_1P_2} = (1000, 0, 0)$,
$\overrightarrow{P_1P_3} = (0, 500, -100)$, $\vec n = (0, 100\,000,
500\,000)$, $\cos\delta = 500\,000/509\,902 \approx 0.981$, $\delta \approx
\ang{11.3}$, et $(n_x, n_y) = (0, 100\,000)$ pointe vers le nord~: même résultat.
\end{solution}

\begin{solution}{exo:b1:geological-maps:6}
L'aval est le nord, $s = 0.03$ (\ang{1.72}). (a) $a = 0 < s$~: V vers l'amont
(sud), le long des courbes de niveau. (b) $a = -\tan\ang{20} < 0 < s$~: V vers
l'amont. (c) $a = \tan\ang{1} \approx 0.017 < s$~: V vers l'amont. (d) $a = \tan\ang{20}
\approx 0.364 > s$~: V vers l'aval (nord).
\end{solution}

\begin{solution}{exo:b1:geological-maps:7}
$\xi_{\text{fond}} - \xi_{\text{versant}} = h/(a - s) = 80/(0.364 - 0.03)
\approx \qty{240}{m}$~: la trace sur le versant est à environ \qty{240}{m} en
amont du sommet du V.
\end{solution}

\begin{solution}{exo:b1:geological-maps:8}
La fermeture de l'\omterm{def:b1:geological-maps:fold}{anticlinal} pointe dans le sens du plongement~: l'axe plonge
vers le nord. Avec un axe plongeant vers le nord, les \omterm{def:b1:geological-maps:geological-map}{affleurements} d'un
\omterm{def:b1:geological-maps:fold}{synclinal} se referment au sud.
\end{solution}

\begin{solution}{exo:b1:geological-maps:9}
Le côté est-sud-est montre des unités plus jeunes (craie et argilite en face
d'unités plus anciennes)~: ce bloc est descendu. Les tirets du symbole de faille
sont de ce côté, le côté affaissé.
\end{solution}

\begin{solution}{exo:b1:geological-maps:10}
Une droite commune aux deux plans est perpendiculaire aux deux normales, donc
parallèle à $\vec n_1 \times \vec n_2$. Pour $090/30$~: $\vec n_1 = (0.577, 0,
1)$~; pour $270/30$~: $\vec n_2 = (-0.577, 0, 1)$. Alors $\vec n_1 \times \vec n_2
= (0\cdot 1 - 1\cdot 0,\ 1\cdot(-0.577) - 0.577\cdot 1,\ 0) = (0, -1.155,
0)$~: une droite horizontale nord--sud. Le premier plan plonge vers l'est, le
second vers l'ouest~: ils plongent en s'éloignant de la droite, qui est la crête
d'un \omterm{def:b1:geological-maps:fold}{anticlinal} à axe horizontal.
\end{solution}

\begin{solution}{exo:b1:geological-maps:11}
La trace du dyke vertical est la droite passant par l'origine, de \omterm{def:b1:geological-maps:attitude}{direction}
\ang{045}, quelle que soit l'altitude~: sur la colline, elle se trouve quelque
part sur cette même droite. Le dyke incliné est le plan $z = -\tan\ang{60}\,d$,
où $d$ est la distance horizontale à la ligne de \omterm{def:b1:geological-maps:attitude}{direction} passant par
l'origine, comptée positivement vers le sud-est. Pour $z = \qty{300}{m}$~:
$d = -300/\tan\ang{60}
\approx \qty{-173}{m}$~: la trace se trouve à \qty{173}{m} au nord-ouest de la
ligne de \omterm{def:b1:geological-maps:attitude}{direction} passant par l'origine, du côté opposé au \omterm{def:b1:geological-maps:attitude}{sens du pendage}.
\end{solution}

\begin{solution}{exo:b1:geological-maps:12}
\[
  \frac{\partial z}{\partial x} \approx \frac{405 - 420}{100} = -0.15, \qquad
  \frac{\partial z}{\partial y} \approx \frac{412 - 420}{100} = -0.08,
\]
donc la pente est $\arctan\sqrt{0.15^2 + 0.08^2} = \arctan 0.17 \approx
\ang{9.6}$. Le sol
descend le long de $-\nabla z = (0.15, 0.08)$, dont l'\omterm{def:b1:spherical-astronomy:horizontal}{azimut} est
$\arctan(0.15/0.08) \approx \ang{62}$~: la maille est tournée vers
l'est-nord-est.
\end{solution}

\begin{solution}{pb:b1:geological-maps:1}
\textbf{1.} Schiste, granite, grès, calcaire, argilite, craie, alluvions.

\textbf{2.} D'ouest en est~: craie et argilite, calcaire, grès, puis schiste et
granite près de $x = \qty{3.5}{km}$, grès, calcaire et argilite de nouveau, puis
craie~; des alluvions sur tout le fond.

\textbf{3.} Le grès, avec le schiste et le granite sous lui le long de la
rivière~: les unités les plus anciennes au cœur, un \omterm{def:b1:geological-maps:fold}{anticlinal}.

\textbf{4.} La craie, l'unité la plus jeune~: un \omterm{def:b1:geological-maps:fold}{synclinal}.

\textbf{5.} La base du grès est le plus haut sur l'axe de l'\omterm{def:b1:geological-maps:fold}{anticlinal}, et ce
n'est que le long de la rivière que l'érosion est descendue sous elle.

\textbf{6.} Une discordance~: le grès repose sur deux unités plus anciennes
différentes, le schiste et le granite, et tronque le contour du granite.

\textbf{7.} $435 - 261 = 174$, $432 - 428 = 4$, $512 - 586 = \qty{-74}{m}$.

\textbf{8.} $\overrightarrow{E_1E_2} = (600, 1200, -170)$,
$\overrightarrow{E_1E_3} = (1100, 400, -248)$.

\textbf{9.} Le produit vectoriel est
\[
  (-229\,600,\ -38\,200,\ -1\,080\,000);
\]
orienté vers le haut, $\vec n_E = (229\,600,\ 38\,200,\ 1\,080\,000)$.

\textbf{10.} $\|\vec n_E\| \approx 1\,104\,800$; $\cos\delta \approx 0.9776$,
$\delta \approx \ang{12.2}$.

\textbf{11.} \omterm{def:b1:spherical-astronomy:horizontal}{Azimut} de $(229\,600, 38\,200)$~: $\arctan(229\,600/38\,200)
\approx \ang{80.6}$~; \omterm{def:b1:geological-maps:attitude}{attitude} $081/12$, \omterm{def:b1:geological-maps:attitude}{direction} $\approx \ang{351}$.

\textbf{12.} Oui~: le symbole du flanc est a un trait de \omterm{def:b1:geological-maps:attitude}{direction} proche de
nord--sud et un tiret de \omterm{def:b1:geological-maps:attitude}{pendage} vers l'est, accompagné de 12.

\textbf{13.} $524 - 669 = -145$, $421 - 459 = -38$, $319 - 209 =
\qty{110}{m}$.

\textbf{14.} $\overrightarrow{W_1W_2} = (700, 1200, 107)$ et
$\overrightarrow{W_1W_3} = (1100, -600, 255)$. Leur produit vectoriel, orienté
vers le haut, est
\[
  (370\,200,\ -60\,800,\ -1\,740\,000) \quad\longrightarrow\quad
  \vec n_W = (-370\,200,\ 60\,800,\ 1\,740\,000).
\]

\textbf{15.} $\cos\delta = 1\,740\,000/\|\vec n_W\| \approx 0.9776$, $\delta
\approx \ang{12.2}$~; \omterm{def:b1:geological-maps:attitude}{sens du pendage}~: \omterm{def:b1:spherical-astronomy:horizontal}{azimut} de $(-370\,200, 60\,800)$,
$\approx \ang{279.3}$. \omterm{def:b1:geological-maps:attitude}{Attitude} $279/12$, \omterm{def:b1:geological-maps:attitude}{direction} $\approx \ang{189}$~: le
flanc ouest plonge vers l'ouest-nord-ouest, le flanc est vers l'est-nord-est.

\textbf{16.} $\vec n_E \times \vec n_W = (804 \times 10^6,\ -799\,320 \times
10^6,\ 28\,101 \times 10^6)$ (arrondi).

\textbf{17.} La \omterm{def:b1:geological-maps:attitude}{direction} vers le bas est $(-804, 799\,320, -28\,101) \times
10^6$~: \omterm{def:b1:geological-maps:attitude}{direction} $\arctan(-804/799\,320) \approx \ang{0}$, c'est-à-dire le nord~;
plongement
$\arctan\bigl(28\,101/\sqrt{804^2 + 799\,320^2}\bigr) \approx \ang{2.0}$. L'axe
plonge de \ang{2} vers le nord.

\textbf{18.} $8000 \times \tan\ang{2.0} \approx \qty{280}{m}$ plus bas, près de
\qty{-50}{m}. Le cœur de l'\omterm{def:b1:geological-maps:fold}{anticlinal} est le plus haut près de la rivière et
s'enfonce vers le nord, si bien que son \omterm{def:b1:geological-maps:geological-map}{affleurement} se ferme au nord, comme le
montre la carte.

\textbf{19.} $\overrightarrow{F_1F_2} = (1627, 2500, -276)$ et
$\overrightarrow{F_1F_3} = (3124, 5500, 77)$ donnent la normale dirigée vers le haut
\[
  (1\,710\,500,\ -987\,503,\ 1\,138\,500), \qquad
  \delta = \arccos\frac{1\,138\,500}{2\,279\,700} \approx \ang{60},
\]
et un \omterm{def:b1:geological-maps:attitude}{sens du pendage} $\approx \ang{120}$. \omterm{def:b1:geological-maps:attitude}{Attitude} $120/60$, \omterm{def:b1:geological-maps:attitude}{direction}
\ang{030}.

\textbf{20.} La faille plonge vers l'est-sud-est, donc ce côté est le
compartiment supérieur, le bloc situé au-dessus du plan de faille incliné~; il
est descendu (unités plus jeunes en surface)~: une faille normale.

\textbf{21.} Le plan de faille plonge vers l'est-sud-est, vers l'aval et bien
plus fortement que ne descend le fond de la rivière~: le terrain plus bas
rencontre le plan plus loin dans le \omterm{def:b1:geological-maps:attitude}{sens du pendage}, si bien que la trace oblique
vers l'est-sud-est dans la vallée --- la règle des V appliquée au plan de faille.

\textbf{22.} Intrusion du granite~; érosion avant le grès~; dépôt de la craie~;
plissement~; jeu de la faille~; dépôt des alluvions.

\textbf{23.} Après 90 millions d'années (la craie est plissée) et avant
2.58 millions d'années (les alluvions ne sont pas faillées)~; le plissement a eu
lieu d'abord, le jeu de la faille ensuite.

\textbf{24.} La faille descend de \qty{1405}{m} de F1 (\qty{405}{m}) à
\qty{-1000}{m}~; avec un \omterm{def:b1:geological-maps:attitude}{pendage} de \ang{60}, cela représente une distance
horizontale de $1405/\tan\ang{60} \approx \qty{811}{m}$ dans le \omterm{def:b1:geological-maps:attitude}{sens du pendage}.
Le puits doit être foré à environ \textbf{\qty[text-series-to-math=true]{810}{m} à l'est-sud-est de F1} (\omterm{def:b1:spherical-astronomy:horizontal}{azimut} \ang{120}).
\end{solution}
